\emph{autoAI} es una aplicación web que acerca la consulta de un catálogo técnico
de vehículos a usuarios sin conocimientos de automoción: permite expresar las
necesidades en lenguaje natural y devuelve recomendaciones fundamentadas en datos
reales, no en afirmaciones inventadas por el modelo de lenguaje.

El sistema se organiza en cuatro servicios desacoplados y contenerizados: un
proceso de \emph{web scraping} con Scrapy sobre \texttt{km77.com}, un
\emph{backend} en FastAPI que normaliza los datos y orquesta el asistente, una
base de datos PostgreSQL y un \emph{frontend} en React y TypeScript. El núcleo es
un \emph{pipeline} basado en el patrón de uso de herramientas (\emph{tool
calling}): el modelo interpreta la petición y redacta la respuesta, pero toda
cifra procede obligatoriamente de la base de datos, lo que impide por
construcción las alucinaciones.

La conclusión principal es que la flexibilidad de un modelo de lenguaje y el
rigor factual de una base de datos pueden conciliarse cuando el modelo se limita
a comprender y redactar y todo dato se delega en la fuente estructurada.
